% ============================================================
%  RESUMEN
% ============================================================
\chapter*{Resumen}
\addcontentsline{toc}{chapter}{Resumen}
\thispagestyle{plain}

Las minas antipersona~(MAP) constituyen una de las amenazas humanitarias más
persistentes en Colombia, país que registra más de 12\,673 víctimas entre 1990 y 2026~\cite{aicma2026}.
Las técnicas de detección convencionales ---detectores de metales, radares de
penetración terrestre y perros entrenados--- presentan limitaciones críticas frente
a minas de fabricación artesanal con contenido metálico mínimo o nulo, además de
exponer al personal operativo a riesgos directos. Este trabajo de grado desarrolla
un sistema de visión artificial basado en imágenes multiespectrales adquiridas
a \SI{1}{m} de altura, para la detección remota de \ac{MAP} tipo quiebrapatas
enterradas superficialmente en suelo arenoso con vegetación limitada.

El sistema utiliza una cámara MicaSense RedEdge-MX de cinco bandas espectrales
(\SI{475}{nm}, \SI{560}{nm}, \SI{668}{nm}, \SI{717}{nm} y \SI{842}{nm}). Antes de
cada sesión de campo se captura un panel de reflectancia calibrada (CRP) que permite
convertir en el procesamiento posterior, los valores digitales de cada banda en
valores de reflectancia de superficie independientes de las condiciones de iluminación.
La base de datos comprende seis sesiones de campo (S2--S7, 14\,901 ventanas)
adquiridas en distintas condiciones ambientales, con validación
\textit{Leave-One-Session-Out}~(LOSO) para evitar fuga de datos entre sesiones; un
experimento complementario que deja fuera cada emplazamiento físico
(\textit{leave-one-zone-out}) confirma que el modelo generaliza a minas en sitios no vistos.
Se emplean cuatro métricas: el \ac{AUC} ---área bajo la curva ROC--- como métrica
principal, por medir la capacidad discriminativa de forma independiente del umbral
y permitir comparar arquitecturas en igualdad de condiciones; y, a umbral fijo~(0.5),
la sensibilidad (fracción de minas detectadas), la especificidad (fracción de
no-minas correctas) y los falsos positivos por mina, como métricas operativas.
Debido al desbalance 30:1, el entrenamiento emplea submuestreo aleatorio~1:1
(EasyEnsemble promedia diez sub-clasificadores balanceados 1:1), lo que permite
reportar correctamente esas métricas operativas a umbral~0.5.

El etiquetado del groundtruth se realizó por anotación manual con bounding box en
Roboflow: una ventana es positiva solo si su centro cae dentro del recuadro de la
\ac{MAP}, lo que produjo 487 ventanas positivas frente a 14\,414 negativas
(desbalance 30:1). El análisis de importancia mostró que las bandas Red Edge y \ac{NIR} aportan a
la discriminación; la selección final de características se realizó con la ganancia
del propio XGBoost, evitando la inconsistencia de usar las importancias de un
clasificador para alimentar otro. Tras comparar \textit{Random Forest},
\ac{SVM}, XGBoost y EasyEnsemble, el modelo final ---\textbf{EasyEnsemble~x10
XGBoost} con las \textbf{70 características} de mayor ganancia--- alcanza
AUC~=~$0.947 \pm 0.007$, con sensibilidad del \textbf{82.3~\%} a umbral~0.5,
comparable con los trabajos previos del grupo \ac{PSI} en termografía infrarroja
(82.4~\%). El análisis por profundidad muestra que la sensibilidad disminuye al
aumentar la profundidad (del 87.5~\% a \SI{1}{cm} al 70.7~\% a \SI{7}{cm}),
coherente con que la señal proviene de la perturbación superficial. El análisis de
falsos positivos muestra que la botella de vidrio es el objeto que más confusión
genera, aunque el \textbf{93~\%} de los falsos positivos en zonas de mina no se
asocian a ningún objeto de interferencia, sino al suelo perturbado alrededor de la mina.

\vspace{6pt}
\noindent\textbf{Palabras clave:} visión artificial; minas antipersona; imágenes
multiespectrales; Red Edge; infrarrojo cercano; estrés vegetal; XGBoost; bosque
aleatorio; validación LOSO.

\clearpage

% ============================================================
%  ABSTRACT
% ============================================================
\chapter*{Abstract}
\addcontentsline{toc}{chapter}{Abstract}
\thispagestyle{plain}

Anti-personnel landmines~(APL) remain one of the most persistent humanitarian threats
in Colombia, a country that has recorded more than 12,673 victims between 1990 and 2026~\cite{aicma2026}
. Conventional detection techniques ---metal detectors, ground-penetrating radars,
and trained dogs--- exhibit critical limitations when applied to artisanal mines with
minimal or no metallic content, while also exposing field personnel to direct risks.
This thesis develops a computer vision system based on multispectral imagery acquired
at \SI{1}{m} height for the remote detection of shallow-buried legbreaker APL in sandy
soil with limited vegetation.

The system uses a MicaSense RedEdge-MX five-band camera (\SI{475}{nm}, \SI{560}{nm},
\SI{668}{nm}, \SI{717}{nm}, and \SI{842}{nm}). Before each field session, a Calibrated
Reflectance Panel~(CRP) is captured to allow the subsequent processing to convert the
digital values of each band into surface reflectance values independent of illumination
conditions. The official dataset comprises six field sessions (S2--S7, 14,901 patches)
acquired under varying environmental conditions, with Leave-One-Session-Out~(LOSO)
validation to prevent data leakage between sessions; a complementary leave-one-zone-out
experiment confirms that the model generalizes to mines at unseen sites. Four metrics are used: the \ac{AUC} ---area under the ROC curve--- as the primary
metric, for measuring discriminative capacity independently of the threshold and
enabling fair comparison of architectures; and, at a fixed threshold~(0.5),
sensitivity (fraction of mines detected), specificity (fraction of non-mines
correctly classified) and false positives per mine, as operational metrics. Due to
the 30:1 imbalance, training uses random 1:1 undersampling (EasyEnsemble averages
ten 1:1-balanced sub-classifiers), which allows these operational metrics to be
reported correctly at the 0.5 threshold.

Groundtruth labeling was performed by manual bounding-box annotation in Roboflow:
a window is positive only if its center falls inside the \ac{MAP} box, yielding 487
positive windows against 14,414 negatives (30:1 imbalance). The importance analysis showed that the Red Edge and Near-Infrared~(NIR) bands
contribute to discrimination; final feature selection used XGBoost's own gain importance,
avoiding the inconsistency of using one classifier's importances to feed another. After comparing
\textit{Random Forest}, \ac{SVM}, XGBoost, and EasyEnsemble, the final model
---\textbf{EasyEnsemble~x10 XGBoost} with the \textbf{top-70 features} by gain---
achieves AUC~=~$0.947 \pm 0.007$, with \textbf{82.3~\%} sensitivity at a 0.5
threshold, comparable with previous work by the PSI research group using infrared
thermography (82.4~\%). Depth analysis shows that sensitivity decreases with burial
depth (from 87.5~\% at 1~cm to 70.7~\% at 7~cm), consistent with a signal
originating from the surface perturbation. The false-positive analysis shows that
the glass bottle is the object generating the most confusion, although \textbf{93~\%}
of false positives in mine zones are not associated with any interference object but with
the disturbed soil around the mine.

\vspace{6pt}
\noindent\textbf{Keywords:} computer vision; anti-personnel landmines; multispectral
imaging; Red Edge; near-infrared; vegetation stress; XGBoost; random forest;
LOSO validation.

\clearpage
